\documentclass[../DoAn.tex]{subfiles}
\begin{document}
{%
\hyphenpenalty=10000
\exhyphenpenalty=10000

\begin{center}
    \Large{\textbf{ABSTRACT}}\\
\end{center}
\vspace{1cm}
In underwater inspection, observation, and manipulation tasks, Remotely Operated Vehicles (ROVs) are commonly controlled manually based on visual feedback from onboard cameras. However, underwater environments pose several challenges, including poor lighting, image noise, water currents, visual latency, and the unavailability of GPS-based positioning. These factors make target tracking and target approach difficult, especially when precise relative motion is required. Existing approaches include fully manual control, distance-sensor-based control, deep-learning-based object detection, and traditional computer vision methods for motion tracking. While deep learning methods can provide strong recognition capability, they often require large training datasets and considerable computational resources. Manual control, on the other hand, depends heavily on the operator’s experience and may become unstable under latency or disturbance. In this graduation project, an Optical Flow-based tracking approach combined with PID control is selected because it is lightweight, suitable for real-time operation, and compatible with the existing ROV platform. The proposed system consists of three main modules: an MJPEG camera streaming module on the ROV, a host-side tracking and target selection module, and an onboard ROV control module. After the operator selects a target on the web interface, the tracking module estimates the target’s displacement and scale variation across consecutive frames. These values are normalized into vision signals, including horizontal displacement, scale error, and tracking confidence, and are transmitted to the controller via UDP. The controller then uses PID algorithms to generate forward/backward and lateral motion commands, while also integrating heading hold and depth hold functions. The final thruster commands are computed through a mixing matrix and sent to the ROV’s thrusters. The main contribution of this project is the development of a complete Track and Catch pipeline that integrates real-time computer vision with ROV motion control. The implemented system enables target selection, visual tracking, confidence-based validation, command limitation, manual override, and target-relative position holding under experimental conditions.

}% end hyphenpenalty group
\end{document}